\section{Introduction}
\label{sec:intro}

How does one trust a surrogate solver for a stochastic PDE? The
standard answer --- validate against a numerical reference --- carries
a failure mode that is easy to state and, as it turns out, easy to
measure. In the companion paper~\cite{paper1} we ran two independent,
individually reasonable discretisations (ADI with upwind advection;
MUSCL--minmod TVD finite volume) of the same white-noise-driven
nonlinear SPDE on the \emph{identical} noise realisation: they differ
by $0.69$ relative $L^2$ pathwise and by a factor $2.1$ in field
variance at the working grid, and for space--time white noise in two
dimensions there is no continuum limit to arbitrate between them ---
the target is lattice-defined, and the lattice is solver-flavoured.
Every reference-validated claim inherits that flavour silently. The
structural remedy is an \emph{a posteriori certificate}: a computable
bound on the distance to the solution of the \emph{declared} operator
equation, with explicit constants, that never routes trust through a
second solver.

\paragraph{This paper.} We construct such a certificate, together with
the method it certifies, for the linear bandlimited-Laplacian SPDE
\begin{equation}
\partial_t u = \Delta u + \sigma\,\dot W_{K_{\max}},
\qquad u(\cdot, 0) = u_0,
\label{eq:lin-spde}
\end{equation}
on the torus $\mathbb{T}^d \times [0, T]$, where $\dot W_{K_{\max}}$ is
a spatially bandlimited Gaussian white noise. In the Whittle--Mat\'ern
spectral representation
$u(x, t) = \sum_{k \in \mathcal{H}_{K_{\max}}}
[a_k(t)\cos(k \cdot x) + b_k(t)\sin(k \cdot x)]$, each coefficient
$a_k, b_k$ is an Ornstein--Uhlenbeck process, and the SPDE decouples
into $J = |\mathcal{H}_{K_{\max}}|$ independent scalar problems in
time. Working in this basis is the design decision from which
everything follows: the spatial Laplacian becomes the diagonal
multiplication by $-|k|^2$, the Galerkin system is an explicit
$L \times L$ solve per mode, and --- because the certificate is
assembled per mode --- \emph{the spatial dimension enters only through
the mode count}. Dimension-independence of the certified bound is not
an empirical surprise to be discovered but a property of the
architecture to be verified; our experiments verify it across a
$444\times$ growth of the mode set. The price of the construction is
the choice of a time basis $\{\psi_l(t)\}$ in which the per-mode
equation is tested --- and, as we show, that choice can fail silently.

\paragraph{Context.} The method sits at the meeting point of two
literatures. Residual-variational physics-informed networks (RVPINNs)
minimise the residual in a discrete dual norm against a finite test
space and yield reliable, efficient a posteriori estimators
\cite{RVPINN2024,Rojas2023}, of the form
$\|u_\theta - u_{\rm ref}\|_V \le C\,\|R(u_\theta)\|_{V_h^*}$ with $C$
an inf-sup constant; that literature concentrates on deterministic
PDEs in $d \le 3$. Spectral (mode-space) representations for SPDE
surrogates were developed in the distributional-regulariser line
\cite{paper1,paper5,paper6}, whose theoretical guarantee is a
Mishra--Molinaro-style \cite{Mishra2022} law-distance ($W_2$) bound
with no path-level control; spectral-domain networks have also been
proposed for deterministic high-dimensional PDEs \cite{YuOseledets2026},
without per-instance error control. A priori error bounds for
physics-informed and operator learning are given by De~Ryck and Mishra
\cite{DeRyckMishra2022}; the present certificate is the complementary
a posteriori object --- evaluated on the computed field rather than
guaranteeing that a good one exists --- and its constant is the
computable, sharp realisation of their abstract stability constant.
For the linear SPDE treated here no network and no training are
involved at all: the certified object is a closed-form Galerkin solve,
and the neural parameterisations of the companion papers appear only
in the nonlinear extension of Section~\ref{sec:discussion}.

\subsection{The time basis can fail silently: plateau and cure}
\label{sec:intro:halfperiod}

A natural choice for the time basis is the standard Dirichlet sine
basis $\psi_l(t) = \sqrt{2/T}\sin(\pi l t/T)$, $l = 1, \ldots, L$,
which vanishes at both endpoints. With this basis the trial-space
ansatz
$a_{\theta,k}(t) = a_{0,k} e^{-|k|^2 t} + \sum_l \alpha_{k,l}\psi_l(t)$
satisfies $a_{\theta,k}(T) = a_{0,k} e^{-|k|^2 T}$ exactly, by
construction; it cannot represent the SPDE's terminal noise
accumulation $x_k(T) = \sigma\int_0^T e^{-|k|^2(T-s)}\,dB(s)$. The
Galerkin formulation in this basis produces a linear system that
admits a unique solution, but the resulting coefficients $\alpha$
\emph{differ from the $L^2$-projection coefficients} of the true OU
process by a tail-correction term that does not vanish as $L$ grows.
Empirically, the relative $L^2(0,T;H^1)$ error stalls at a plateau of
$\approx 0.33$ regardless of $L$, up to $L = 1024$
(Section~\ref{sec:exp:dirichlet-plateau}) --- a manifestation of the
non-self-adjointness of the parabolic operator
$\partial_t + |k|^2\,\mathrm{I}$ under this restricted test space.
No amount of compute repairs it, and without an error instrument the
failure is invisible: the Galerkin system remains well-posed and its
residual small in the wrong norm. The certificate is precisely the
instrument that exposes it.

The cure is architectural: replace the Dirichlet sine basis with the
\emph{half-period sine basis}
\begin{equation}
\phi_l(t) = \sqrt{2/T}\,\sin\!\left(\frac{(2l-1)\pi t}{2T}\right),
\qquad l = 1, 2, \ldots, L,
\label{eq:half-period-basis}
\end{equation}
which satisfies $\phi_l(0) = 0$ but $\phi_l(T) = (-1)^{l-1}$ --- free
at the terminal time. These are the eigenfunctions of $-\partial_t^2$
with Dirichlet at $t = 0$ and Neumann at $t = T$, the natural boundary
conditions for the parabolic problem. With matched boundary conditions
the Galerkin error attains the \emph{optimal} Karhunen--Lo\`eve rate
for the rough (H\"older-$\tfrac12$) OU reference,
$\mathrm{rel\_err} \approx 1.6/\sqrt{L}$, uniformly across
$d \in \{2, 3, 5, 7\}$ (Section~\ref{sec:exp:d-sweep}). The
corresponding certified bound is
\begin{equation}
\|u_\theta - u_{\rm ref}\|_Y^2
\;\le\; C_1\,\mathcal{L}_{\rm RV}(\theta) + \frac{C_2}{L},
\label{eq:main-bound}
\end{equation}
with $C_1$ from the discrete inf-sup constant of the matched Galerkin
formulation and $C_2$ from the best-approximation error of the OU
process in the half-period basis. Both constants are uniform in $d$.

\subsection{Contributions}
\label{sec:intro:contrib}

\begin{enumerate}
\item \textbf{A certified mode-space Galerkin method, uniform in $d$
by construction} (Sections~\ref{sec:method}, \ref{sec:theorem}).
Trial/test bases that make the Gram matrix block-diagonal in $k$ and
the per-mode time operator an explicit $L \times L$ system; the
certificate $\|u_\theta - u_{\rm ref}\|_Y^2 \le
C_1\mathcal{L}_{\rm RV} + C_2/L$ assembles per mode, so $d$ enters
only through the mode count. Closed-form solve; no autograd, no FFT,
no training.

\item \textbf{A silent structural failure, diagnosed and cured}
(Section~\ref{sec:exp:dirichlet-plateau}). The natural Dirichlet
time basis stalls at a non-decaying plateau ($\approx 0.33$ up to
$L = 1024$); the root cause is identified at the coefficient level
(Galerkin $\ne$ $L^2$ projection, non-vanishing tail correction);
the half-period basis removes the obstruction and attains the
optimal Karhunen--Lo\`eve rate $\approx 1.6/\sqrt{L}$ for the rough
reference.

\item \textbf{Verification of the designed dimension-independence,
synthetic and real} (Sections~\ref{sec:exp:d-sweep},
\ref{sec:exp:realworld}). Across $d \in \{2, 3, 5, 7\}$ ($J$ growing
$444\times$) the relative error is unchanged within $8\%$ at every
$L$; on ERA5 temperature and SPX implied-volatility data the same
uniformity holds, with the absolute rate set by the data's roughness
and noise floor.
\end{enumerate}

\subsection{Roadmap}
\label{sec:intro:roadmap}

Section~\ref{sec:method} describes the mode-space architecture, the
half-period basis, and the closed-form linear solve. Section
\ref{sec:theorem} states and proves Theorem~\ref{thm:main}. Section
\ref{sec:experiments} reports the empirical validation
(\ref{sec:exp:dirichlet-plateau} the negative result;
\ref{sec:exp:d-sweep} the uniform-in-$d$ verification). Section
\ref{sec:discussion} discusses extensions to nonlinear SPDEs.
